% \section{INTRODUCTION}
% \label{sec:introduction}

% % \subsection{Background and Motivation}
% The demand for privacy-preserving collaborative learning has increased rapidly as NLU systems are deployed across mobile devices, enterprise platforms, and edge services. In many practical settings, NLU data contains sensitive user information (e.g., personal messages, medical text, behavioral logs), making centralized data collection difficult due to privacy policies and regulatory constraints. FL addresses this issue by keeping raw data on client devices and sharing only model updates \cite{mcmahan2017fedavg}.

% At the same time, distributed AI systems are now expected to support many users, domains, and tasks in parallel \cite{mcmahan2017fedavg,smith2017mocha}. This trend makes centralized single-pipeline training less suitable and motivates federated training frameworks that can learn collaboratively while respecting data locality.

% \subsection{Challenges in Federated Natural Language Understanding}
% Although FL is promising for NLU, three major challenges remain.

% First, client data is typically non-IID. Different users and organizations produce language data with distinct vocabularies, label distributions, and styles, which often slows convergence and reduces global model quality \cite{li2020fedprox,smith2017mocha}.

% Second, task heterogeneity is common in realistic deployments. A federation may contain intent classification, sentiment analysis, question answering, or generation-related objectives simultaneously \cite{smith2017mocha,chen2024fedbone,zhou2025mfed}. Classical one-model FL methods are not designed for this mixed-task regime.

% Third, communication and scalability constraints are fundamental. Frequent synchronization across many clients can create substantial overhead, especially when model sizes are large and client connectivity is unstable \cite{mcmahan2017fedavg,chen2024fedbone}.

% Recent federated MTL literature (e.g., MOCHA, FedBone, MIRA, and encoder--decoder based M-Fed) highlights that both data heterogeneity and task heterogeneity must be treated as first-class design factors \cite{smith2017mocha,chen2024fedbone,elbakary2025mira,zhou2025mfed}.

% % \subsection{Role of Small Language Models in Federated Systems}
% Small Language Models (SLMs) are especially attractive in federated systems because they improve system efficiency without removing the benefits of language pretraining.

% Key advantages include:
% \begin{itemize}
%     \item Communication efficiency: fewer parameters reduce upload/download payload per communication round.
%     \item Lower computation cost: SLMs require less local memory and training time, enabling broader client participation.
%     \item Edge-device deployability: compact models are easier to run and update on resource-constrained devices.
% \end{itemize}

% In practice, this role extends to parameter-efficient adaptation techniques (e.g., adapter-style and low-rank updates), which preserve competitive performance while reducing federated overhead.

% \subsection{Real-Time Federated Learning for NLU}
% Most current FL pipelines are batch-oriented: clients train for multiple local steps and synchronize periodically \cite{mcmahan2017fedavg,li2020fedprox}. This design is effective for static datasets but less suited to streaming and dynamic NLU environments where language patterns, topics, and user intent drift over time.

% Real-time federated learning aims to shorten feedback loops by supporting faster synchronization, timely adaptation, and more responsive model updates. For online NLU applications (e.g., conversational assistants, live moderation, event-driven analytics), delayed global updates can quickly degrade performance. Therefore, real-time synchronization mechanisms are increasingly important for maintaining both relevance and robustness in production systems.

% \subsection{Contributions of This Work}
% Building on the above motivation and the reviewed literature, this work provides:
% \begin{itemize}
%     \item A real-time FL-MTL framework for heterogeneous NLU tasks.
%     \item Integration of small language models in federated environments to improve efficiency and deployability.
%     \item Analysis of how data heterogeneity affects FL-MTL convergence and downstream task performance.
%     \item Evaluation across multiple NLU benchmark tasks to assess local and global effectiveness.
% \end{itemize}

% The remainder of this paper is organized as follows. Section~\ref{sec:related_work} reviews relevant literature in federated and MTL. Section~\ref{sec:problem_definition} provides the formal definition of the RT-FMTL problem. Section~\ref{sec:proposed_method} details the system architecture and optimization strategies. Section~\ref{sec:experiments} presents the experimental setup and analysis of results. Finally, Section~\ref{sec:conclusion} concludes the paper and discusses future work.

\section{INTRODUCTION}
\label{sec:introduction}

With the proliferation of NLU systems on mobile devices, enterprise platforms, and edge services, the demand for privacy-preserving collaborative learning has grown significantly \cite{li2020challenges}. Centralized collection of NLU data from sources such as personal messages, medical text, or behavioral logs is often restricted by strict privacy regulations and user concerns. Consequently, FL has emerged as a key solution by keeping raw data local and sharing only model updates \cite{mcmahan2017fedavg}. Modern NLU applications, ranging from sentiment analysis for customer feedback to real-time intent classification for conversational agents and semantic similarity for information retrieval, increasingly require the ability to learn from diverse, distributed data streams while supporting MTL simultaneously \cite{smith2017mocha,chen2024fedbone}. The trend makes centralized single-pipeline training less feasible and encourages FL training models, which can learn collaboratively without breaking the data locality.

While FL is a promising technique for NLU, there are three major concerns, as highlighted by benchmarking frameworks like FedNLP \cite{lin2021fednlp}. First, client data usually is non-IID. Given that diverse users and organizations write language data with different vocabularies, label compositions and preferences, convergence is typically slowed, and global model quality is poor \cite{li2020fedprox,smith2017mocha}. Second, task heterogeneity is common in realistic deployments. A federation could be made with intent classification, sentiment analysis, question answering, or even generation-related goals all at once \cite{smith2017mocha,chen2024fedbone,zhou2025mfed}. Classical one-model FL methods are not suited to this mixed-task regime. Third, communication and scalability are essential. In general, the repeated syncing of multiple clients can lead to significant overhead, particularly when the size of the model is large and the connected clients unstable \cite{mcmahan2017fedavg,chen2024fedbone}. Recent literature on federated MTL (e.g., MOCHA \cite{smith2017mocha}, FedBone \cite{chen2024fedbone}, MIRA \cite{elbakary2025mira}, and encoder--decoder based M-Fed \cite{zhou2025mfed}) indicates that data heterogeneity and task heterogeneity need to be considered as first-class design considerations.

SLMs are especially appealing in federated systems because they improve efficiency without losing language pretraining functionality. Key advantages include:
\begin{itemize}
    \item Communication efficiency: A smaller parameter set reduces the communication cost per round (both upload and download).
    \item Reduced computation cost: SLMs use fewer resources of local memory and training time, allowing for broader client involvement.
    \item Edge-device deployment: compact models are easier to deploy and update on resource-constrained devices.
\end{itemize}
This role translates into parameter-efficient adaptation techniques (e.g., adapter-style and low-rank updates) that preserve competitive performance while decreasing federated overhead.

Current FL pipelines predominantly rely on batch-based synchronization, where clients execute multiple local steps and synchronize periodically \cite{mcmahan2017fedavg,li2020fedprox}. While robust for static datasets, this design is less suited for streaming and dynamic environments, where language patterns, subjects, and user intent evolve over time. To address these limitations, this paper proposes RT-FedMTL, a real-time FL-MTL framework specifically designed for low-latency synchronization and agile adaptation. By utilizing WebSocket protocols and optimizing for SLMs, RT-FedMTL reduces the feedback loop, enabling near-instantaneous global updates that are essential for responsive applications like chatbots, live moderation, and event-driven analytics.

This work makes several key contributions to the field of federated NLU:
\begin{itemize}
    \item We introduce RT-FedMTL, a novel real-time federated framework that integrates MTL and SLMs to leverage cross-task knowledge transfer while minimizing on-device resource footprints.
    \item We evaluate five SLM architectures (DistilBERT \cite{sanh2019distilbert}, BERT-Medium \cite{devlin2018bert}, BERT-Mini \cite{devlin2018bert}, MiniLM \cite{wang2020minilm}, and TinyBERT \cite{jiao2019tinybert}) across three benchmark NLU tasks (sentiment analysis, paraphrase detection, and semantic similarity) to identify the optimal balance between performance and resource consumption.
    \item We analyze the impact of non-IID and task heterogeneity on convergence speed and task accuracy in a real-time environment.
    \item We demonstrate that RT-FedMTL achieves a significant reduction in per-client GPU memory usage (up to 74--93\%) compared to centralized MTL baselines while maintaining high performance across multiple benchmarks from the GLUE suite.
\end{itemize}

The rest of the paper is structured as follows. Section~\ref{sec:related_work} reviews the relevant work in FL and MTL. Section~\ref{sec:problem_definition} defines the formal problem of RT-FedMTL. In Section~\ref{sec:proposed_method} we detail the system architecture and optimization strategies. The experimental setup and results analysis are described in Section~\ref{sec:experiments}. Finally, Section~\ref{sec:conclusion} wraps up the paper and presents the next steps.